% =====================================================================
%  MANUAL DE USUARIO — ATARA
%  Sistema de Alerta Temprana y Atención al Rendimiento Académico
%  Dirigido a: Docentes de Educación Básica — MEP Costa Rica
% =====================================================================
\documentclass[12pt,a4paper,oneside]{article}

% --------------------- Codificación e idioma ---------------------------
\usepackage[utf8]{inputenc}
\usepackage[T1]{fontenc}
\usepackage[spanish,es-tabla]{babel}
\usepackage{lmodern}
\usepackage{microtype}

% --------------------- Geometría y estilo ------------------------------
\usepackage[a4paper,top=2.8cm,bottom=2.8cm,left=3.0cm,right=2.5cm,
            headheight=25pt]{geometry}
\usepackage{fancyhdr}
\usepackage{parskip}
\setlength{\parskip}{6pt}

% --------------------- Títulos -----------------------------------------
\usepackage[explicit]{titlesec}
\usepackage{titletoc}

% --------------------- Tablas y listas ---------------------------------
\usepackage{booktabs}
\usepackage{array}
\usepackage{tabularx}
\usepackage{enumitem}
\usepackage{multirow}

% --------------------- Colores, cajas, imágenes ------------------------
\usepackage{xcolor}
\usepackage[most]{tcolorbox}
\usepackage{graphicx}
\usepackage{float}
\usepackage{tikz}
\usepackage{pifont}

% --------------------- Paleta de colores ATARA -------------------------
\definecolor{ataraAzul}{HTML}{1F3A93}
\definecolor{ataraAzulClaro}{HTML}{4A6FBF}
\definecolor{ataraAzulPale}{HTML}{EEF2FB}
\definecolor{ataraVerde}{HTML}{1B7F3B}
\definecolor{ataraVerdePale}{HTML}{EAF5EE}
\definecolor{ataraNaranja}{HTML}{D97E1F}
\definecolor{ataraNaranjaPale}{HTML}{FDF3E3}
\definecolor{ataraRojo}{HTML}{B5292F}
\definecolor{ataraRojoPale}{HTML}{FCEAEA}
\definecolor{ataraGris}{HTML}{6B7280}
\definecolor{ataraGrisPale}{HTML}{F3F4F6}
\definecolor{ataraAmarillo}{HTML}{F59E0B}

% --------------------- Hipervínculos -----------------------------------
\usepackage[colorlinks=true,linkcolor=ataraAzul,urlcolor=ataraAzul]{hyperref}

% --------------------- Estilo de títulos de sección -------------------
\titleformat{\section}
  {\color{ataraAzul}\Large\bfseries}{}{0em}
  {\colorbox{ataraAzul}{\color{white}\makebox[\dimexpr\linewidth-2\fboxsep][l]
   {\enspace\thesection\enspace\textbar\enspace #1\enspace}}}
\titlespacing{\section}{0pt}{16pt}{8pt}

\titleformat{\subsection}
  {\color{ataraAzul}\large\bfseries}{\thesubsection.}{0.5em}{#1}
  [\color{ataraAzulClaro}\rule{\linewidth}{0.5pt}]
\titlespacing{\subsection}{0pt}{12pt}{4pt}

\titleformat{\subsubsection}
  {\color{ataraAzulClaro}\normalsize\bfseries}{\thesubsubsection.}{0.5em}{#1}
\titlespacing{\subsubsection}{0pt}{8pt}{2pt}

% --------------------- Encabezado y pie --------------------------------
\pagestyle{fancy}
\fancyhf{}
\fancyhead[L]{\small\color{ataraAzul}\textbf{ATARA} \textcolor{ataraGris}{— Manual de Usuario}}
\fancyhead[R]{\small\color{ataraGris}MEP Costa Rica}
\fancyfoot[C]{\small\thepage}
\renewcommand{\headrulewidth}{0.4pt}
\renewcommand{\headrule}{%
  \hbox to\headwidth{\color{ataraAzul}\leaders\hrule height\headrulewidth\hfill}}

% --------------------- Cajas de contenido -----------------------------
\newtcolorbox{consejo}[1][]{%
  colback=ataraVerdePale, colframe=ataraVerde!60,
  arc=3mm, boxrule=0.6pt, left=6pt, right=6pt, top=4pt, bottom=4pt,
  fonttitle=\bfseries\small\color{ataraVerde},
  title={{\large\ding{51}}\enspace Consejo}, #1}

\newtcolorbox{importante}[1][]{%
  colback=ataraNaranjaPale, colframe=ataraNaranja!70,
  arc=3mm, boxrule=0.6pt, left=6pt, right=6pt, top=4pt, bottom=4pt,
  fonttitle=\bfseries\small\color{ataraNaranja},
  title={{\large\textbf{!}}\enspace Importante}, #1}

\newtcolorbox{aviso}[1][]{%
  colback=ataraRojoPale, colframe=ataraRojo!60,
  arc=3mm, boxrule=0.6pt, left=6pt, right=6pt, top=4pt, bottom=4pt,
  fonttitle=\bfseries\small\color{ataraRojo},
  title={{\large\ding{55}}\enspace Atención}, #1}

\newtcolorbox{nota}[1][]{%
  colback=ataraAzulPale, colframe=ataraAzul!50,
  arc=3mm, boxrule=0.6pt, left=6pt, right=6pt, top=4pt, bottom=4pt,
  fonttitle=\bfseries\small\color{ataraAzul},
  title={{\large\textbf{i}}\enspace Nota}, #1}

\newtcolorbox{pasos}[1][Pasos]{%
  colback=ataraGrisPale, colframe=ataraAzulClaro!70,
  arc=3mm, boxrule=0.8pt, left=8pt, right=6pt, top=4pt, bottom=6pt,
  fonttitle=\bfseries\color{ataraAzul},
  title=#1}

% --------------------- Placeholder para imágenes ----------------------
\newcommand{\imagen}[2][8cm]{%
  \begin{center}
  \begin{tikzpicture}
    \draw[fill=ataraGrisPale, draw=ataraAzulClaro!60, dashed, line width=1pt,
          rounded corners=4pt]
      (0,0) rectangle (0.92\linewidth,#1);
    \node[align=center, text=ataraGris, text width=0.82\linewidth] at (current bounding box.center) {%
      \large\itshape #2};
  \end{tikzpicture}
  \end{center}}

\newcommand{\imagenS}[1]{%
  \begin{center}
  \begin{tikzpicture}
    \draw[fill=ataraGrisPale, draw=ataraAzulClaro!60, dashed, line width=1pt,
          rounded corners=4pt]
      (0,0) rectangle (0.75\linewidth,5.5cm);
    \node[align=center, text=ataraGris, text width=0.65\linewidth] at (current bounding box.center) {%
      \normalsize\itshape #1};
  \end{tikzpicture}
  \end{center}}

% --------------------- Comandos utilitarios ----------------------------
\newcommand{\boton}[1]{\fcolorbox{ataraAzulClaro}{ataraAzulPale}%
  {\small\textbf{\color{ataraAzul} #1}}}
\newcommand{\menu}[1]{\textbf{\color{ataraAzulClaro} #1}}
\newcommand{\campo}[1]{\textit{«#1»}}

\newenvironment{pasosnum}{%
  \begin{enumerate}[leftmargin=1.6cm, label=\textbf{\color{ataraAzul}\arabic*.},
                    itemsep=4pt, parsep=0pt]}
  {\end{enumerate}}

% =====================================================================
%                              PORTADA
% =====================================================================
\begin{document}
\pagenumbering{roman}
\thispagestyle{empty}

\begin{tikzpicture}[remember picture, overlay]
  \fill[ataraAzul] (current page.north west) rectangle
    ([yshift=-7.5cm]current page.north east);
\end{tikzpicture}

\vspace*{1.8cm}

\begin{center}
  {\Huge\bfseries\color{white} ATARA}\\[0.5em]
  {\large\color{white!85!ataraAzul}
    Sistema de Alerta Temprana y Atención al Rendimiento Académico}\\[0.3em]
  {\normalsize\color{white!70!ataraAzul}
    División de Educación Básica — MEP Costa Rica}

  \vspace{2.5cm}

  {\Huge\bfseries\color{ataraAzul} Manual de Usuario}\\[0.6em]
  {\Large\color{ataraGris} Guía para Docentes}

  \vspace{1.5cm}

  \begin{tcolorbox}[
    width=0.80\linewidth,
    colback=ataraAzulPale,
    colframe=ataraAzul!40,
    arc=4mm, boxrule=0.8pt,
    halign=center]
  {\normalsize Este manual le explica paso a paso cómo usar ATARA
  para registrar evaluaciones por saber, generar alertas temáticas
  y dar seguimiento al rendimiento de sus estudiantes.}
  \end{tcolorbox}

  \vspace{2.5cm}

  \begin{tabular}{rl}
  \textbf{Versión:}          & 1.1 \\[0.3em]
  \textbf{Fecha:}            & Mayo 2026 \\[0.3em]
  \textbf{Dirigido a:}       & Docentes de Educación Básica \\[0.3em]
  \textbf{Acceso al sistema:}& \texttt{https://atara-backend.vercel.app} \\
  \end{tabular}
\end{center}

\newpage

% =====================================================================
%                              ÍNDICE
% =====================================================================
\tableofcontents
\newpage
\pagenumbering{arabic}

% =====================================================================
\section{Bienvenida y descripción general}
% =====================================================================

Bienvenido a \textbf{ATARA}. Esta plataforma fue creada especialmente para que usted,
como docente, pueda llevar un seguimiento detallado del aprendizaje de sus estudiantes
y detectar de manera oportuna aquellos que necesitan apoyo adicional.

\subsection{¿Para qué sirve ATARA?}

Con ATARA usted puede:

\begin{itemize}[leftmargin=1.5cm, itemsep=4pt]
  \item[\textbf{1.}] \textbf{Consultar sus secciones asignadas} — ver el listado de grupos
        y estudiantes que tiene a su cargo durante el año lectivo.
  \item[\textbf{2.}] \textbf{Registrar evaluaciones por saber} — valorar a cada estudiante
        en los ejes temáticos de las materias que imparte (escala del 1 al 5).
  \item[\textbf{3.}] \textbf{Ver alertas temáticas automáticas} — el sistema calcula
        promedios y le avisa cuáles ejes presentan riesgo de bajo rendimiento.
  \item[\textbf{4.}] \textbf{Consultar visualizaciones y reportes} — ver gráficas del
        rendimiento de su sección y obtener reportes detallados por estudiante.
\end{itemize}

\subsection{Roles en el sistema}

ATARA distingue tres tipos de usuarios. Como docente, su acceso está \textbf{limitado a las
secciones y materias que le fueron asignadas} por el administrador.

\begin{table}[H]
\centering\renewcommand{\arraystretch}{1.3}\small
\begin{tabularx}{\textwidth}{@{}p{3.2cm} X@{}}
\toprule
\textbf{Rol} & \textbf{¿Qué puede hacer?} \\
\midrule
\textbf{Docente}
  & Consultar sus secciones, evaluar estudiantes, ver alertas,
    visualizaciones y reportes de sus grupos asignados. \\
\textbf{Coordinador}
  & Las mismas funciones que el docente. \\
\textbf{Administrador}
  & Acceso total al sistema: gestión de usuarios, años lectivos,
    secciones, estudiantes e importación de listas PIAD. \\
\bottomrule
\end{tabularx}
\end{table}

\begin{nota}
Si no puede ver una sección o materia que debería poder ver, comuníquese con el
administrador del sistema para que le asigne el acceso correspondiente.
\end{nota}

% =====================================================================
\section{Acceso al sistema}
% =====================================================================

\subsection{Cómo ingresar}

\begin{pasosnum}
  \item Abra su navegador web y vaya a la dirección:
        \texttt{https://atara-backend.vercel.app}
  \item Ingrese su \campo{correo institucional} y \campo{contraseña}.
  \item Haga clic en \boton{Ingresar}.
\end{pasosnum}

\imagen[5.5cm]{[Captura de pantalla: Pantalla de inicio de sesión]}

Al ingresar, el sistema lo llevará directamente a la sección de
\textbf{Secciones}, donde verá los grupos que tiene asignados.

\begin{consejo}
Si es la primera vez que ingresa, utilice las credenciales que le proporcionó el
administrador del sistema. Si tiene problemas para acceder, contáctelo directamente.
\end{consejo}

\subsection{Cómo cerrar sesión}

Para salir del sistema de forma segura, haga clic en el \textbf{ícono con su inicial}
ubicado en la esquina superior derecha de la pantalla. Se desplegará un menú con la
opción \boton{Cerrar sesión}.

\begin{aviso}
No cierre el navegador directamente sin cerrar sesión, especialmente en computadoras de uso
compartido.
\end{aviso}

\subsection{Navegación general}

El sistema tiene un \textbf{menú lateral} en el lado izquierdo de la pantalla con acceso
rápido a todas las secciones disponibles. Como docente verá únicamente las opciones
que le corresponden.

\imagen[6cm]{[Captura de pantalla: Menú lateral con las opciones del docente]}

\begin{table}[H]
\centering\renewcommand{\arraystretch}{1.3}\small
\begin{tabularx}{\textwidth}{@{}p{4.5cm} X@{}}
\toprule
\textbf{Opción del menú} & \textbf{¿Para qué sirve?} \\
\midrule
\menu{Secciones}           & Ver sus grupos asignados y los estudiantes de cada uno \\
\menu{Eval.\ por Saber}    & Registrar y consultar evaluaciones de sus estudiantes \\
\menu{Alertas Tempranas}   & Ver alertas generadas por bajo rendimiento \\
\menu{Visualizaciones}     & Ver gráficas del rendimiento por sección \\
\menu{Reportes}            & Obtener reportes individuales por estudiante \\
\menu{Mi sesión}           & Consultar la información de su cuenta activa \\
\bottomrule
\end{tabularx}
\end{table}

% =====================================================================
\section{Mis Secciones}
% =====================================================================

La pantalla de \textbf{Secciones} es el punto de partida para el trabajo diario. Aquí
verá todos los grupos de estudiantes que le han sido asignados para el año lectivo activo.

\subsection{Cómo ver sus secciones}

\begin{pasosnum}
  \item En el menú lateral, haga clic en \menu{Secciones}.
  \item Verá la lista de grupos asignados a usted (por ejemplo: «4°A», «5°B»).
  \item Haga clic sobre una sección para ver la lista de estudiantes matriculados.
\end{pasosnum}

\imagenS{[Captura de pantalla: Lista de secciones asignadas al docente]}

\begin{nota}
Si una sección que imparte no aparece en su lista, comuníquese con el administrador
del sistema para que le asigne ese grupo a su cuenta.
\end{nota}

% =====================================================================
\section{Evaluaciones por Saber \textnormal{\small(Función principal)}}
\label{sec:evalsaber}
% =====================================================================

Esta es la funcionalidad más importante de ATARA para el docente. Aquí usted registra,
periodo a periodo, cómo va cada estudiante en cada \emph{eje temático} de la materia
que imparte.

\subsection{¿Qué es una evaluación por saber?}

Cada evaluación corresponde a \textbf{un estudiante}, \textbf{una materia} y \textbf{un
periodo lectivo}. Dentro de ella, usted asigna una valoración del \textbf{1 al 5} a cada
eje temático según el tipo de saber que esté evaluando.

\begin{table}[H]
\centering\renewcommand{\arraystretch}{1.25}\small
\begin{tabular}{@{}c l p{9cm}@{}}
\toprule
\textbf{Valor} & \textbf{Nivel} & \textbf{Significa\ldots} \\
\midrule
\textbf{1} & Inicial       & El estudiante no evidencia dominio del eje \\
\textbf{2} & En desarrollo & Dominio parcial, necesita apoyo constante \\
\textbf{3} & Intermedio    & Dominio básico, comete algunos errores \\
\textbf{4} & Logrado       & Buen dominio, errores mínimos \\
\textbf{5} & Avanzado      & Dominio pleno y puede transferir el conocimiento \\
\bottomrule
\end{tabular}
\caption{Escala de valoración utilizada en las evaluaciones por saber.}
\end{table}

\subsection{Tipos de saber}

Cada materia tiene ejes temáticos organizados en tres categorías:

\begin{itemize}[leftmargin=1.5cm, itemsep=3pt]
  \item \textbf{Saber Conceptual} — lo que el estudiante \emph{sabe} (conocimientos teóricos).
  \item \textbf{Saber Procedimental} — lo que el estudiante \emph{sabe hacer} (habilidades prácticas).
  \item \textbf{Saber Actitudinal} — cómo el estudiante \emph{actúa y se relaciona} con el aprendizaje.
\end{itemize}

Puede evaluar a un mismo estudiante en los tres tipos de saber dentro del mismo periodo.
Cada tipo genera una evaluación independiente.

\subsection{Registrar una nueva evaluación}

\begin{pasos}[Cómo registrar una evaluación por saber]
\begin{pasosnum}
  \item En el menú lateral, haga clic en \menu{Eval.\ por Saber}.
  \item Seleccione la \campo{sección} de los estudiantes que desea evaluar.
  \item Seleccione el \campo{periodo lectivo} activo.
  \item Haga clic en \boton{+ Nueva evaluación}.
  \item Seleccione el \campo{estudiante} y la \campo{materia}.
  \item Elija el \campo{tipo de saber} (Conceptual, Procedimental o Actitudinal).
  \item Asigne una valoración del 1 al 5 a cada eje temático mostrado.
  \item Si todos los ejes tienen el mismo valor, use \boton{Aplicar nivel a todos los ejes}
        para ahorrar tiempo.
  \item Haga clic en \boton{Guardar}.
\end{pasosnum}
\end{pasos}

\imagen[7cm]{[Captura de pantalla: Formulario de nueva evaluación por saber — selector de sección y período]}

\imagen[7cm]{[Captura de pantalla: Wizard de evaluación — asignación de valores por eje temático]}

\begin{importante}
Si cierra el formulario con datos ingresados sin guardar, el sistema le pedirá
confirmación antes de descartar los cambios. Esto evita pérdidas accidentales de información.
\end{importante}

\subsection{Consultar evaluaciones existentes}

\begin{pasosnum}
  \item En \menu{Eval.\ por Saber}, seleccione la sección y el periodo.
  \item La lista mostrará todas las evaluaciones registradas para ese grupo.
  \item Haga clic sobre una evaluación para ver sus detalles completos.
\end{pasosnum}

\imagen[6cm]{[Captura de pantalla: Lista de evaluaciones por saber de una sección]}

\subsection{Modificar una evaluación (recalificación)}

Si necesita corregir los valores de una evaluación ya guardada:

\begin{pasosnum}
  \item Localice la evaluación en la lista y haga clic en \boton{Editar}.
  \item Modifique los valores de los ejes que desea corregir.
  \item Haga clic en \boton{Guardar cambios}.
\end{pasosnum}

\begin{nota}
Al editar una evaluación existente se reemplazan los valores anteriores. El sistema
\textbf{no} crea una segunda evaluación; actualiza la existente. El historial de
cambios queda registrado internamente.
\end{nota}

% =====================================================================
\section{Alertas Temáticas \textnormal{\small(Función principal)}}
% =====================================================================

Las alertas temáticas son la herramienta más poderosa de ATARA. Una vez que usted ha
registrado evaluaciones por saber, el sistema calcula automáticamente qué estudiantes
presentan \textbf{riesgo de bajo rendimiento} en cada eje temático y le genera un aviso.

\subsection{¿Qué es una alerta temática?}

El sistema calcula el \textbf{promedio} de todas las valoraciones de un estudiante en
cada eje temático y determina un nivel de alerta según ese promedio:

\begin{table}[H]
\centering\renewcommand{\arraystretch}{1.35}
\begin{tabular}{@{}c p{3.4cm} p{7.2cm}@{}}
\toprule
\textbf{Promedio} & \textbf{Nivel} & \textbf{¿Qué significa?} \\
\midrule
$\leq 2{,}0$ &
  \colorbox{ataraRojoPale}{\textbf{\color{ataraRojo}\quad ALTA \quad}} &
  Riesgo crítico. Este estudiante requiere atención inmediata en ese eje. \\[0.4em]
$2{,}1 - 3{,}0$ &
  \colorbox{ataraNaranjaPale}{\textbf{\color{ataraNaranja}\quad MEDIA \quad}} &
  Riesgo moderado. Se recomienda seguimiento cercano. \\[0.4em]
$> 3{,}0$ &
  \colorbox{ataraVerdePale}{\textbf{\color{ataraVerde}\ SIN ALERTA\ }} &
  Desempeño aceptable. No se genera alerta para este eje. \\
\bottomrule
\end{tabular}
\caption{Niveles de alerta temática y su interpretación pedagógica.}
\end{table}

\subsection{Generar alertas para un estudiante}

\begin{pasos}[Cómo generar alertas para un estudiante]
\begin{pasosnum}
  \item En el menú lateral, haga clic en \menu{Alertas Tempranas}.
  \item Seleccione la \campo{sección} y el \campo{periodo}.
  \item Haga clic en el nombre del estudiante.
  \item Haga clic en \boton{Generar alertas}.
  \item El sistema analizará todas las evaluaciones registradas y mostrará las alertas
        correspondientes por eje temático.
\end{pasosnum}
\end{pasos}

\imagen[6.5cm]{[Captura de pantalla: Vista de Alertas Temáticas con lista de estudiantes]}

\subsection{Generar alertas para toda la sección}

Si desea analizar a todos los estudiantes de una sección a la vez:

\begin{pasosnum}
  \item Seleccione la \campo{sección} y el \campo{periodo}.
  \item Haga clic en \boton{Generar alertas para la sección}.
  \item El sistema procesará a cada estudiante matriculado automáticamente.
\end{pasosnum}

\begin{consejo}
Genere las alertas al finalizar cada trimestre o cuando haya registrado varias
evaluaciones nuevas. El proceso puede repetirse las veces que sea necesario; las
alertas anteriores se actualizan con los datos más recientes.
\end{consejo}

\subsection{Consultar e imprimir las alertas}

\begin{pasosnum}
  \item En \menu{Alertas Tempranas}, seleccione la sección y el periodo.
  \item Verá la lista de estudiantes con sus alertas activas, ordenadas por nivel de riesgo.
  \item Para imprimir o guardar como PDF, haga clic en \boton{Imprimir}.
\end{pasosnum}

\imagen[7cm]{[Captura de pantalla: Lista de alertas temáticas por estudiante con colores por nivel (ALTA/MEDIA)]}

\begin{importante}
Los ejes con nivel \textbf{SIN ALERTA} no aparecen en la lista — el sistema muestra
únicamente los casos que requieren su atención (ALTA y MEDIA), para que pueda
enfocarse en lo que realmente importa.
\end{importante}

% =====================================================================
\section{Visualizaciones}
% =====================================================================

Las visualizaciones le muestran de un vistazo cómo está rindiendo su sección
completa en una materia y periodo determinados. Es ideal para identificar tendencias
del grupo antes de planificar sus próximas clases.

\subsection{Cómo acceder a las visualizaciones}

\begin{pasosnum}
  \item Haga clic en \menu{Visualizaciones} en el menú lateral.
  \item Seleccione la \campo{sección}, la \campo{materia} y el \campo{periodo}.
  \item El sistema mostrará una gráfica con la distribución de niveles de desempeño
        de todos los estudiantes del grupo.
\end{pasosnum}

\imagen[7cm]{[Captura de pantalla: Gráfica de distribución de desempeño de una sección]}

\subsection{¿Cómo interpretar la gráfica?}

La gráfica muestra cuántos estudiantes de su sección se encuentran en cada nivel de
desempeño (Inicial, En desarrollo, Intermedio, Logrado, Avanzado) para la materia
seleccionada. Esto le permite identificar rápidamente si la mayor parte del grupo
está progresando o si hay una concentración en los niveles más bajos que requiere
atención.

\begin{consejo}
Use las visualizaciones al final de cada trimestre para preparar su informe de
avance del grupo y planificar estrategias de apoyo diferenciadas para el siguiente periodo.
\end{consejo}

% =====================================================================
\section{Reportes por estudiante}
% =====================================================================

Los reportes le permiten obtener un resumen completo del desempeño de un
estudiante específico: sus promedios por eje temático, sus alertas activas y el
historial de evaluaciones del periodo seleccionado.

\subsection{Cómo generar un reporte}

\begin{pasosnum}
  \item Haga clic en \menu{Reportes} en el menú lateral.
  \item Busque y seleccione el \campo{estudiante}.
  \item Seleccione la \campo{materia} y el \campo{periodo}.
  \item Haga clic en \boton{Generar reporte}.
  \item El reporte se mostrará en pantalla. Puede imprimirlo con
        \boton{Imprimir} o guardarlo como PDF desde su navegador
        (Archivo → Imprimir → Guardar como PDF).
\end{pasosnum}

\imagen[6.5cm]{[Captura de pantalla: Reporte individual de un estudiante con promedios y alertas]}

\begin{nota}
El reporte incluye: promedios por tipo de saber, detalle de cada eje temático evaluado
y las alertas activas del periodo seleccionado. Es especialmente útil para reuniones con
madres y padres de familia o para el expediente del estudiante.
\end{nota}

% =====================================================================
\section{Alertas Tempranas \textnormal{\small(Módulo complementario)}}
% =====================================================================

La sección de \textbf{Alertas Tempranas} es un módulo complementario que permite
consultar un historial general de seguimiento académico por estudiante o sección,
basado en una escala de cuatro puntos.

\subsection{Escala del módulo de Alertas Tempranas}

\begin{table}[H]
\centering\renewcommand{\arraystretch}{1.2}\small
\begin{tabular}{@{}c l@{}}
\toprule
\textbf{Valor} & \textbf{Nivel} \\
\midrule
1 & Insuficiente \\
2 & Básico \\
3 & Satisfactorio \\
4 & Destacado \\
\bottomrule
\end{tabular}
\end{table}

\subsection{Consultar alertas tempranas}

\begin{pasosnum}
  \item Haga clic en \menu{Alertas Tempranas} en el menú lateral.
  \item Seleccione el \campo{estudiante} o la \campo{sección}.
  \item Seleccione el \campo{periodo} de interés.
  \item El sistema mostrará las alertas generadas con este modelo de evaluación.
  \item Para imprimir, haga clic en \boton{Imprimir}.
\end{pasosnum}

\imagenS{[Captura de pantalla: Vista de Alertas Tempranas de un estudiante]}

\begin{nota}
Para el registro de nuevas evaluaciones, utilice preferentemente el módulo
\textbf{Evaluaciones por Saber} (sección \ref{sec:evalsaber} de este manual),
que corresponde al modelo pedagógico actualizado del MEP con la escala del 1 al 5.
\end{nota}

% =====================================================================
\section{Importar lista PIAD \textnormal{\small(Solo administradores)}}
% =====================================================================

El módulo de importación PIAD permite cargar el PDF oficial de la lista de
adecuaciones del MEP (\emph{Plan Individual de Adecuación}) para extraer
automáticamente la lista de estudiantes con adecuaciones curriculares.

\begin{importante}
Esta función está disponible únicamente para usuarios con rol de
\textbf{Administrador}. Si usted es docente y necesita registrar estudiantes con
adecuaciones, comuníquese con el administrador del sistema.
\end{importante}

\subsection{¿Qué es el PIAD?}

El PIAD es el documento oficial del MEP que lista a los estudiantes que tienen
adecuaciones curriculares. ATARA puede leer ese PDF y extraer los nombres y cédulas
de forma automática mediante reconocimiento de texto (OCR).

\subsection{Cómo importar la lista PIAD}

\begin{pasos}[Cómo cargar el PDF de la lista PIAD]
\begin{pasosnum}
  \item En el menú lateral, haga clic en \menu{Importar PIAD}.
  \item Haga clic en \boton{Seleccionar archivo} y elija el PDF de la lista PIAD.
  \item Haga clic en \boton{Extraer estudiantes}.
  \item El sistema mostrará la lista de estudiantes encontrados en el PDF.
  \item Revise la lista para confirmar que los datos son correctos.
  \item Registre manualmente en \menu{Estudiantes} los que no figuren aún en el sistema.
\end{pasosnum}
\end{pasos}

\imagen[6.5cm]{[Captura de pantalla: Módulo de importación PIAD con área de carga del PDF]}

\imagen[5.5cm]{[Captura de pantalla: Lista de estudiantes extraídos del PDF]}

\begin{aviso}
El archivo debe ser el PDF oficial generado por el sistema SIGCE del MEP. PDFs
editados o generados por otros programas pueden no ser reconocidos correctamente.
\end{aviso}

% =====================================================================
\section{Gestión de estudiantes \textnormal{\small(Solo administradores)}}
% =====================================================================

En esta sección el administrador puede buscar, registrar y actualizar la información
básica de los estudiantes del sistema.

\begin{importante}
Esta función está disponible únicamente para usuarios con rol de \textbf{Administrador}.
Los docentes consultan a sus estudiantes directamente desde la sección \menu{Secciones}.
\end{importante}

\subsection{Buscar un estudiante}

\begin{pasosnum}
  \item Haga clic en \menu{Estudiantes} en el menú lateral.
  \item Utilice el buscador para encontrar al estudiante por nombre o cédula.
  \item Haga clic en el nombre del estudiante para ver su información.
\end{pasosnum}

\imagenS{[Captura de pantalla: Lista de estudiantes con buscador]}

\subsection{Registrar un estudiante manualmente}

\begin{pasosnum}
  \item En \menu{Estudiantes}, haga clic en \boton{+ Nuevo estudiante}.
  \item Complete los campos: nombre, apellidos, número de cédula y fecha de nacimiento.
  \item Haga clic en \boton{Guardar}.
\end{pasosnum}

\begin{nota}
Si un estudiante ya existe en el sistema con la misma cédula, \textbf{no se creará un
duplicado}. El sistema lo reconoce y devuelve el registro existente automáticamente.
\end{nota}

% =====================================================================
\section{Preguntas frecuentes}
% =====================================================================

\begin{description}[leftmargin=1cm, style=nextline, font=\bfseries\color{ataraAzul},
                    itemsep=10pt]

  \item[¿Qué hago si no veo mi sección en el sistema?]
  Comuníquese con el administrador del sistema. Es posible que su cuenta aún no
  tenga asignada esa sección o que el año lectivo no esté activo.

  \item[¿Puedo evaluar a un estudiante en más de una materia?]
  Sí, pero únicamente en las materias que le han sido asignadas por el administrador.
  Si necesita acceso a una materia adicional, solicítelo directamente.

  \item[¿Qué pasa si ingreso un valor equivocado en una evaluación?]
  Puede editarla en cualquier momento. Vaya a \menu{Eval.\ por Saber}, busque la
  evaluación correspondiente y haga clic en \boton{Editar}.

  \item[¿Con qué frecuencia debo generar las alertas temáticas?]
  Se recomienda generarlas al finalizar cada trimestre o cada vez que haya
  registrado evaluaciones nuevas para un grupo. El proceso es rápido y puede
  repetirse sin problema; los resultados se actualizan automáticamente.

  \item[¿Puedo imprimir las alertas para llevarlas a una reunión?]
  Sí. Tanto en \menu{Alertas Tempranas} como en \menu{Reportes} encontrará un
  botón \boton{Imprimir} que genera una vista lista para impresora o para guardar
  como PDF desde el navegador.

  \item[No veo la opción «Importar PIAD» ni «Estudiantes» en mi menú. ¿Por qué?]
  Estas opciones están disponibles únicamente para administradores del sistema.
  Como docente, sus estudiantes ya están registrados por el administrador; usted
  los consulta desde \menu{Secciones}.

  \item[Olvidé mi contraseña. ¿Qué hago?]
  Comuníquese con el administrador del sistema para que le restablezca el acceso.
  ATARA no cuenta con recuperación de contraseña automática por correo electrónico
  en esta versión.

  \item[¿Mis evaluaciones se guardan automáticamente?]
  No. Debe hacer clic en \boton{Guardar} para que queden registradas. Si cierra
  el formulario sin guardar, el sistema le preguntará si desea descartar los cambios,
  evitando pérdidas accidentales.

\end{description}

% =====================================================================
\section{Glosario de términos}
% =====================================================================

\begin{description}[leftmargin=3.5cm, style=sameline, font=\bfseries\color{ataraAzul},
                    labelwidth=3.2cm, itemsep=6pt]

  \item[Eje temático]
    Subdivisión de contenido dentro de una materia. Cada materia tiene
    21 ejes temáticos (7 por tipo de saber). Ejemplo: «Álgebra y pensamiento
    algebraico» es un eje de Matemáticas.

  \item[Tipo de saber]
    Categoría del aprendizaje: \emph{Conceptual} (saber), \emph{Procedimental}
    (saber hacer) o \emph{Actitudinal} (saber ser).

  \item[Materia]
    Asignatura que se evalúa. En esta versión: Español, Matemáticas, Ciencias
    y Estudios Sociales.

  \item[Periodo]
    Trimestre lectivo (I, II o III) dentro de un año lectivo.

  \item[Sección]
    Grupo de estudiantes a cargo de un docente. Ejemplo: «4°A».

  \item[Alerta temática]
    Aviso generado automáticamente cuando el promedio de un estudiante en un
    eje temático está en riesgo (ALTA: promedio $\leq 2{,}0$; MEDIA: entre $2{,}1$ y $3{,}0$).

  \item[PIAD]
    Plan Individual de Adecuación. Documento oficial del MEP con la lista de
    estudiantes que tienen adecuaciones curriculares.

  \item[Evaluación por saber]
    Registro del desempeño de un estudiante en los ejes temáticos de una
    materia y periodo, con valoraciones del 1 al 5.

\end{description}

\vspace{1.5cm}

\begin{tcolorbox}[
  colback=ataraAzulPale, colframe=ataraAzul!50,
  arc=4mm, boxrule=0.8pt, halign=center]
{\large\bfseries\color{ataraAzul} ¿Necesita ayuda adicional?}\\[0.4em]
Comuníquese con el coordinador de su institución o con el equipo técnico
responsable del sistema ATARA en su circuito educativo.
\end{tcolorbox}

\end{document}
